% © 2026 Ing. David Jaroš — CC BY-NC-ND 4.0
% Licensed under CC BY-NC-ND 4.0. See LICENSE.md.
\documentclass[12pt,a4paper]{article}
\usepackage[a4paper,margin=2.5cm]{geometry}
\usepackage{amsmath,amssymb,amsthm,mathtools}
\usepackage[most]{tcolorbox}
\usepackage[colorlinks=true,linkcolor=blue]{hyperref}
\newtheorem{definition}{Definition}[section]
\newtheorem{theorem}[definition]{Theorem}
\newtheorem{corollary}[definition]{Corollary}
\newtheorem{proposition}[definition]{Proposition}
\theoremstyle{remark}\newtheorem{remark}[definition]{Remark}
\title{GAP-10$\Omega$: Reconstruction of the Classical Frame Connection}
\author{Ing. David Jaro\v{s}}
\date{16 July 2026}
\begin{document}
\maketitle

\begin{abstract}
The covariant-tetrad formulation of UBT uses
$E_\mu=\mathcal N_0^{-1/2}D_\mu\Theta$ and the central anticommutator metric.
This note proves two reconstruction results.  First, in the nondegenerate
metric-compatible torsion-free classical GR branch, the Lorentz frame
connection is uniquely the Levi--Civita spin connection of the tetrad.  Second,
for arbitrary specified torsion, the complete metric-compatible connection is
still uniquely reconstructed from the tetrad and torsion through the
contorsion tensor.  Hence the connection itself carries no residual kinematic
ambiguity once $(e,T)$ are fixed.  What remains open in full UBT is the
*dynamical law selecting the torsion* and the exact representation by which the
connection acts on $\Theta$.  The flat special-relativistic limit is verified,
and the implicit self-consistency equation created by
$E_\mu=D_\mu\Theta/\sqrt{\mathcal N_0}$ is isolated from the already closed
connection-reconstruction problem.
\end{abstract}

\section{Geometric data}
Let $M$ be a four-dimensional manifold and let
\begin{equation}
 e^a=e_\mu{}^a\,dx^\mu,
 \qquad \det(e_\mu{}^a)\ne0,
 \label{eq:coframe}
\end{equation}
be the real coefficients of the UBT Lorentz-slice tetrad
\begin{equation}
 E_\mu=i e_\mu{}^0\mathbf 1+e_\mu{}^k\mathbf e_k.
\end{equation}
The induced metric is
\begin{equation}
 g_{\mu\nu}=e_\mu{}^a e_\nu{}^b\eta_{ab},
 \qquad \eta_{ab}=\operatorname{diag}(-1,1,1,1).
 \label{eq:metric}
\end{equation}
Let $e_a{}^\mu$ denote the inverse tetrad.

A Lorentz frame connection is a one-form
\begin{equation}
 \omega^a{}_b=\omega_\mu{}^a{}_b\,dx^\mu,
 \qquad \omega_{\mu ab}=-\omega_{\mu ba}.
\end{equation}
Its torsion two-form is
\begin{equation}
 T^a:=de^a+\omega^a{}_b\wedge e^b
      =\frac12T^a{}_{bc}\,e^b\wedge e^c.
 \label{eq:torsion}
\end{equation}
The corresponding biquaternionic/spin connection is the lift
\begin{equation}
 \Omega_\mu=\rho_{\mathrm{spin}*}(\omega_\mu)
 =\frac14\omega_{\mu ab}\,\Sigma^{ab}
 \label{eq:spinlift}
\end{equation}
for the chosen UBT representation.  The normalization of the generators
$\Sigma^{ab}$ is representation-conventional; the Lorentz connection
$\omega_\mu{}^{ab}$ itself is not.

\section{Unique torsion-free compatible connection}
\begin{theorem}[Levi--Civita frame-connection theorem]
\label{thm:unique-omega}
For every $C^2$ nondegenerate tetrad $e_\mu{}^a$, there exists exactly one
Lorentz connection $\mathring\omega_\mu{}^{ab}$ satisfying
\begin{enumerate}
 \item frame-metric compatibility,
       $\mathring\omega_{\mu ab}=-\mathring\omega_{\mu ba}$;
 \item vanishing torsion, $\mathring T^a=0$.
\end{enumerate}
It is given by
\begin{equation}
 \boxed{
 \mathring\omega_\mu{}^a{}_b
 =e_b{}^\nu\left(
 \mathring\Gamma^\rho{}_{\mu\nu}(g)e_\rho{}^a
 -\partial_\mu e_\nu{}^a
 \right),}
 \label{eq:omega-explicit}
\end{equation}
where $\mathring\Gamma^\rho{}_{\mu\nu}(g)$ is the Levi--Civita connection of
Eq.~\eqref{eq:metric}.
\end{theorem}

\begin{proof}
The Levi--Civita connection is the unique affine connection satisfying
$\mathring\nabla_\mu g_{\nu\rho}=0$ and
$\mathring\Gamma^\rho{}_{[\mu\nu]}=0$.  Define $\mathring\omega$ by
Eq.~\eqref{eq:omega-explicit}.  Multiplication by $e_\nu{}^b$ gives the tetrad
compatibility identity
\begin{equation}
 \partial_\mu e_\nu{}^a
 -\mathring\Gamma^\rho{}_{\mu\nu}e_\rho{}^a
 +\mathring\omega_\mu{}^a{}_b e_\nu{}^b=0.
 \label{eq:tetrad-postulate-proof}
\end{equation}
Antisymmetrising in $\mu,\nu$ and using
$\mathring\Gamma^\rho{}_{[\mu\nu]}=0$ gives $\mathring T^a=0$.
Differentiating $g_{\mu\nu}=e_\mu{}^a e_\nu{}^b\eta_{ab}$, using
$\mathring\nabla g=0$ and Eq.~\eqref{eq:tetrad-postulate-proof}, gives
$\mathring\omega_{\mu ab}+\mathring\omega_{\mu ba}=0$.

For uniqueness, let $\omega$ and $\omega'$ satisfy both conditions and define
$K^a{}_b:=\omega^a{}_b-\omega'^a{}_b$.  Metric compatibility gives
$K_{abc}=-K_{bac}$.  Equality of the torsions gives
$K^a{}_b\wedge e^b=0$, hence in the tetrad basis
$K_{abc}=K_{acb}$.  Combining the two symmetries,
\[
 K_{abc}=K_{acb}=-K_{cab}=-K_{cba}=K_{bca}=K_{bac}=-K_{abc},
\]
so $K_{abc}=0$ and therefore $\omega=\omega'$.
\end{proof}

\section{General metric-compatible connection with torsion}
The torsion-free theorem is the special case of a stronger kinematic
reconstruction statement.

\begin{theorem}[Tetrad--torsion reconstruction theorem]
\label{thm:tetrad-torsion}
Fix a $C^2$ nondegenerate tetrad $e_\mu{}^a$ and a torsion tensor
$T^a{}_{bc}=-T^a{}_{cb}$.  There exists exactly one metric-compatible Lorentz
connection with that tetrad and torsion.  It is
\begin{equation}
 \boxed{\omega^a{}_b=\mathring\omega^a{}_b(e)+K^a{}_b(T),}
 \label{eq:omega-decomp}
\end{equation}
where $\mathring\omega(e)$ is Eq.~\eqref{eq:omega-explicit} and, with all frame
indices lowered by $\eta_{ab}$,
\begin{equation}
 \boxed{
 K_{abc}=\frac12\left(T_{cab}-T_{abc}-T_{bca}\right).}
 \label{eq:contorsion-formula}
\end{equation}
The convention is $T^a=de^a+\omega^a{}_b\wedge e^b$ as in
Eq.~\eqref{eq:torsion}.
\end{theorem}

\begin{proof}
Write $\omega=\mathring\omega+K$.  Since $\mathring\omega$ is torsion free,
\begin{equation}
 T^a=K^a{}_b\wedge e^b.
\end{equation}
Metric compatibility of both $\omega$ and $\mathring\omega$ implies
$K_{abc}=-K_{bac}$.  In frame components the torsion relation is equivalently
\begin{equation}
 T_{abc}=K_{acb}-K_{abc}=-(K_{cab}+K_{abc}).
 \label{eq:T-K-cycle}
\end{equation}
Writing the two cyclic permutations of Eq.~\eqref{eq:T-K-cycle} and taking
$T_{cab}-T_{abc}-T_{bca}$ gives $2K_{abc}$, proving
Eq.~\eqref{eq:contorsion-formula}.  Thus existence is explicit.  If two
metric-compatible connections have the same tetrad and torsion, their
contorsions have zero difference by the same formula, proving uniqueness.
\end{proof}

\begin{corollary}[No residual kinematic connection ambiguity]
\label{cor:no-independent-omega}
For a nondegenerate tetrad, a metric-compatible connection is completely fixed
once the torsion is specified.  In particular, in the classical GR branch
$T=0$ implies $K=0$ and $\omega=\mathring\omega(e)$.
\end{corollary}

This result does \emph{not} say that UBT has already derived $T=0$.  It says
that the remaining full-UBT question is the dynamical law for torsion, not an
unconstrained freedom in $\Omega_\mu$ itself.

\section{Gauge covariance and Lorentz-slice preservation}
Under a local Lorentz transformation $e\mapsto\Lambda e$,
\begin{equation}
 \omega\mapsto\Lambda\omega\Lambda^{-1}-d\Lambda\,\Lambda^{-1}.
\end{equation}
This is gauge covariance, not a new physical ambiguity.  Since
$\omega_\mu\in\mathfrak{so}(1,3)$, parallel transport preserves
$\eta_{ab}$.  Therefore the induced vector action preserves the real Lorentz
slice and the central anticommutator metric.  This closes the connection part
of GAP-10L for every metric-compatible branch, including branches with
contorsion.  It does not prove that the full $\Theta$ equations and their
sources dynamically keep $D_\mu\Theta$ in that slice.

\section{Flat limit}
\begin{corollary}[Special relativity]
For the inertial tetrad $e_\mu{}^a=\delta_\mu{}^a$ in Cartesian Minkowski
coordinates,
\begin{equation}
 \mathring\Gamma^\rho{}_{\mu\nu}=0,
 \qquad \mathring\omega_\mu{}^{ab}=0,
 \qquad \Omega_\mu=0,
 \qquad D_\mu=\partial_\mu,
\end{equation}
provided the selected UBT vacuum torsion is zero.  More generally, flat
spacetime may have nonzero connection coefficients in a curvilinear frame,
while its curvature remains zero.
\end{corollary}

\section{Implicit $\Theta$--connection self-consistency}
The reconstruction theorems treat $(e,T)$ as known geometric data.  In UBT,
however,
\begin{equation}
 E_\mu=\mathcal N_0^{-1/2}
 \left(\partial_\mu\Theta+\rho_*(\Omega_\mu[E,T])\Theta\right)
 \label{eq:implicit-self-consistency}
\end{equation}
is an implicit nonlinear first-order system.  Unique reconstruction of
$\Omega_\mu$ does not by itself prove that
Eq.~\eqref{eq:implicit-self-consistency} has a local or global solution for
every required tetrad.

Moreover, the representation $\rho_*$ matters.  A naive one-sided action on an
invertible $\Theta$ creates a flatness obstruction in a torsion-free compatible
branch, whereas a natural two-sided biquaternionic bimodule action can support
nonzero conjugate left/right curvatures.  This is proved separately in
\texttt{canonical/gr\_closure/gap\_10i\_integrability\_selection.tex}.

\begin{tcolorbox}[colback=green!6!white,colframe=green!55!black,title=Defensible status]
\textbf{GAP-10$\Omega$-KIN: CLOSED [L1].}  For every nondegenerate tetrad and
specified torsion, the metric-compatible frame connection is uniquely
$\omega=\mathring\omega(e)+K(T)$, up to local Lorentz gauge.

\medskip
\textbf{GAP-10$\Omega$-GR: CLOSED [L1].}  In the torsion-free classical GR
branch, $K=0$ and the frame connection is uniquely the Levi--Civita spin
connection of the tetrad.

\medskip
\textbf{GAP-10T-DYN: OPEN.}  The canonical UBT action must determine whether
vacuum torsion vanishes and whether matter/spin sources produce a nonzero
contorsion.

\medskip
\textbf{GAP-10L-CONN: CLOSED [L1].}  Any metric-compatible Lorentz connection
preserves $\eta_{ab}$ and the Lorentz slice.  Full dynamical slice preservation
remains GAP-10L-DYN.

\medskip
\textbf{GAP-10I-CURVED: OPEN.}  Existence and on-shell integrability of the
implicit system $E_\mu=D_\mu\Theta/\sqrt{\mathcal N_0}$ for curved geometries
remain to be proved.
\end{tcolorbox}

\end{document}
